
%Table~\ref{tab:TaskATwitterFull} gives a short description of the systems that participated in subtask A, testing on the Twitter dataset;
%we believe this gives a good idea about the approaches, features, lexicons, etc. that the participants used.
%The approaches used for subtask B and for the SMS dataset were similar, and the participating teams also overlapped.
%Thus, here we are taking the opportunity to discuss the characteristics of the systems submitted to subtask A for testing on Twitter data;
%these findings should generalize to subtask B and the SMS dataset.

We have seen that most participants restricted themselves to the provided data and submitted constrained systems.
Indeed, the best systems for each of the two subtasks and for each of the two testing datasets were constrained systems;
of course, this does not mean that additional data would not be useful.
Curiously, in some cases where a team submitted a constrained and unconstrained run,
the unconstrained run actually performed worse.

Not surprisingly, most systems were supervised;
there were only five semi-supervised systems, and there was only one unsupervised system.
One additional team declared their system as unsupervised
since it was not making use of the training data;
we still classified it as supervised though since it did use supervision -- in the form of manual rules.

Most participants predicted all three labels (positive, negative and neutral),
even though some participants opted for not predicting neutral, which made some sense since the final F1-score
was averaged over the positive and the negative predictions only.

The most popular classifiers included SVM, MaxEnt, linear classifier, Naive Bayes;
in some cases, manual rules or ensembles of classifiers were used.

A variety of features were used, including
word-related (e.g., words, stems, $n$-grams, word clusters),
word-shape (e.g., punctuation, capitalization),
syntactic (e.g., POS tags, dependency relations),
Twitter-specific (e.g., repeated characters, emoticons, URLs, hashtags, slang, abbreviations),
and sentiment-related (e.g., negation);
one team also used discourse relations.
Almost all participants relied heavily of various sentiment lexicons,
the most popular ones being MPQA and SentiWordNet,
as well as AFINN and Bing Liu's Opinion Lexicon;
some participants used their own lexicons -- preexisting or built from the provided data.

Given that Twitter messages are noisy, most participants did some preprocessing, including
tokenization,
stemming, lemmatization,
stopword removal,
normalization/removal of URLs, hashtags, users, slang, emoticons, repeated vowels, punctuation;
some even did pronoun resolution.
